\chapter{Introduccion, Problema y Diseno de la Investigacion}

\section{1.1 Contexto e introduccion}
El entrenamiento de redes neuronales profundas ha incrementado de forma sostenida su complejidad computacional, su densidad paramétrica y sus demandas de memoria, hasta el punto de convertir la disponibilidad de recursos hardware en una condición metodológica del propio resultado científico. En la práctica contemporánea del aprendizaje profundo, la GPU concentra la mayor parte de la carga crítica por su capacidad de cómputo paralelo, mientras que la CPU queda frecuentemente relegada a tareas auxiliares. Sin embargo, esta distribución de funciones no es una ley natural del entrenamiento, sino una convención operativa que responde más a la simplicidad de implementación que a una optimización rigurosa del sistema completo.

La consecuencia inmediata de esa convención es bien conocida: la memoria de la GPU se convierte en el cuello de botella dominante durante el entrenamiento de modelos profundos, especialmente cuando se incrementan resolución, longitud de secuencia, número de capas o tamaño del lote. Al mismo tiempo, la CPU y su espacio de memoria principal permanecen infrautilizados como recurso estructural del proceso de entrenamiento. Esta asimetría no solo limita la escalabilidad, sino que introduce una pregunta de investigación legítima: si la carga del entrenamiento puede distribuirse de forma científicamente controlada entre CPU y GPU, entonces la política de colocación deja de ser un detalle de ingeniería y pasa a constituir un objeto formal de estudio. La literatura sobre rematerializacion, checkpointing y particion automatica confirma que esta tension entre memoria, computo y comunicacion es estructural y no contingente \cite{griewank2000revolve,chen2016sublinear,mirhoseini2017device,jia2019beyond}.

Desde una perspectiva abstracta, el entrenamiento puede entenderse como un proceso de minimización de una función de pérdida sobre un espacio paramétrico:

$$
\theta^* = \arg\min_{\theta} \mathcal{L}(\theta; \mathcal{D}),
$$

pero esa formulación, por sí sola, oculta una dimensión material decisiva: la trayectoria hacia $\theta^*$ depende de una realización física concreta del grafo de entrenamiento. La tesis parte precisamente de esta observación. La optimización estadística del modelo y la organización computacional que la hace posible no son problemas independientes, sino dos caras de un mismo proceso. Por ello, la disponibilidad y la asignación de recursos dejan de ser cuestiones puramente instrumentales y pasan a integrarse dentro del razonamiento científico del trabajo.

La presente monografía se sitúa precisamente en ese punto. Su problema rector no es meramente acelerar código ni describir un conjunto de scripts, sino formular una teoría operacional de la partición heterogénea CPU-GPU para entrenamiento profundo bajo restricciones físicas explícitas. El trabajo entiende la heterogeneidad como un espacio de decisión sobre el cual es posible razonar, medir, optimizar y validar. En consecuencia, la tesis adopta una perspectiva de sistemas y de optimización experimental: mide el comportamiento por capa, robustifica coeficientes, construye un modelo declarativo, transforma la solución en un plan ejecutable y contrasta ese plan con evidencia observada.

Desde esta perspectiva, el valor científico del problema reside en que obliga a unir tres dominios que con frecuencia aparecen separados: la fenomenología empírica del hardware, la formalización matemática de la decisión y la verificación operacional de la política resultante. Si cualquiera de estos planos se omite, el resultado queda metodológicamente incompleto. Una optimización sin datos observados se vuelve abstracta; un profiling sin modelo de decisión es descriptivo pero no normativo; y una solución sin ejecución verificable no alcanza estatuto suficiente de evidencia doctoral en una tesis de sistemas.

En consecuencia, la introducción del problema no puede reducirse a una narrativa de escasez de memoria. Lo que está en juego es la posibilidad de construir una teoría aplicada del entrenamiento heterogéneo donde el diseño del sistema deje de ser un conjunto de decisiones implícitas y se convierta en una política formalmente justificable. Esa transición, desde la intuición ingenieril hacia la decisión explícita bajo restricciones, es el punto de partida conceptual de toda la monografía.

Esta formulacion inicial cumple ademas una funcion de encuadre epistemologico. Permite desplazar el foco desde la pregunta instrumental de si una tecnica concreta acelera o no una corrida hacia una pregunta mas fuerte: bajo que condiciones una redistribucion de computo puede considerarse una decision cientificamente defendible. La diferencia no es retorica. Una mejora puntual de rendimiento puede ser valiosa en ingenieria aplicada; una tesis doctoral, en cambio, debe explicar por que esa mejora es interpretable, reproducible y susceptible de generalizacion conceptual.

En este sentido, la introduccion no opera solo como entrada narrativa al manuscrito, sino como delimitacion del tipo de conocimiento que se busca producir. La tesis no pretende ofrecer una receta local de optimizacion para un modelo particular, sino una forma de pensar la ejecucion heterogenea cuando el hardware deja de ser un soporte transparente y se convierte en parte constitutiva del problema. Esta orientacion es la que justifica el paso desde el analisis descriptivo del entrenamiento hacia una arquitectura metodologica que combina observacion, formalizacion y contraste operacional.

\section{1.2 Motivacion tecnica y relevancia cientifica}
La motivación técnica del trabajo nace de una tensión concreta entre dos hechos empíricos. El primero es que el entrenamiento profundo depende de la retención temporal de activaciones, gradientes y estados del optimizador, lo cual produce una presión acumulativa sobre la memoria de la GPU. El segundo es que la CPU coexiste en la misma plataforma con recursos de memoria y cómputo significativos, pero rara vez participa de manera estructurada en la ejecución del grafo de entrenamiento. La tesis parte de la hipótesis de que esta asimetría no expresa un límite arquitectónico inevitable, sino la ausencia de un marco formal que permita explotar la heterogeneidad sin romper la coherencia operacional del entrenamiento.

Ese marco no puede construirse mediante reubicaciones ad hoc de capas ni mediante intuiciones locales basadas solo en throughput. Requiere una visión integrada de cuatro elementos: costos nodales por capa, dependencias estructurales del grafo, penalización de transferencias entre dispositivos y restricciones de memoria sobre ambos subsistemas. En otras palabras, la CPU no puede incorporarse como recurso complementario útil si la tesis no modela simultáneamente latencia, energía, memoria y comunicación. La dificultad técnica del problema es precisamente que toda decisión local de colocación modifica el equilibrio global del sistema. Ese diagnostico tambien aparece, con diferentes lenguajes, en trabajos sobre paralelizacion de pipeline, particion automatica y optimizacion de colocacion, donde la calidad de la decision depende de dependencias globales y no de reglas locales \cite{huang2019gpipe,narayanan2019pipedream,zheng2022alpa}.

Si se denota de manera esquemática el costo total del entrenamiento como una combinación de cómputo, memoria y comunicación,

$$
Z \sim Z_{comp} + Z_{mem} + Z_{comm},
$$

entonces la motivación técnica del trabajo puede formularse como la búsqueda de políticas que no minimicen un único término de forma miope, sino que controlen el compromiso entre los tres. Un desplazamiento agresivo hacia GPU puede reducir $Z_{comp}$ y, a la vez, volver inviable $Z_{mem}$; un reparto excesivo hacia CPU puede aliviar memoria pero inflar $Z_{comm}$ y degradar el régimen temporal del entrenamiento. La tesis adquiere sentido precisamente porque este compromiso no puede resolverse por intuición local, sino mediante un esquema explícito de decisión.

La relevancia científica surge de convertir ese problema práctico en una metodología reproducible y defendible. La tesis no solo pregunta si la CPU puede participar más, sino bajo qué régimen de restricciones, con qué costo de transferencia, con qué efecto sobre la memoria disponible y con qué consecuencias sobre la calidad final del entrenamiento. En este punto el trabajo se aleja deliberadamente de un enfoque de benchmarking convencional: su meta no es producir una tabla de tiempos, sino un marco de inferencia experimental que permita derivar políticas de asignación interpretables y auditables.

En términos más amplios, el proyecto se inscribe en la transición hacia un aprendizaje profundo consciente de recursos. Esta transición supone que la eficiencia ya no puede definirse exclusivamente por latencia o throughput, sino por la capacidad de sostener entrenamiento útil bajo limitaciones materiales reales. Por ello, la tesis se sitúa en la intersección entre sistemas heterogéneos, optimización discreta y metodología experimental, y adopta como criterio de calidad no la mera novedad algorítmica, sino la coherencia entre medición, modelado y validación.

La relevancia científica del tema reside además en su capacidad de generalización conceptual. Aunque el caso de estudio se desarrolla sobre entrenamiento profundo CPU-GPU, la pregunta subyacente es más amplia: cómo convertir restricciones materiales del sistema en variables explícitas de decisión dentro de un marco formal. En este sentido, la tesis no solo aspira a producir una solución técnica utilizable, sino también a proponer un modo de razonar sobre aprendizaje automático cuando el recurso físico deja de ser abundante y pasa a formar parte del problema mismo.

Esa capacidad de generalizacion es particularmente importante en el contexto contemporaneo del aprendizaje profundo. A medida que los modelos crecen, la distancia entre formulacion estadistica y realizacion material se vuelve mas estrecha. Ya no basta con demostrar que un metodo converge o que una arquitectura mejora una metrica final; es necesario mostrar tambien que tal comportamiento puede sostenerse bajo restricciones realistas de memoria, energia y comunicacion. La motivacion cientifica del trabajo se asienta precisamente en esa mutacion del campo: la eficiencia de sistemas deja de ser un apendice ingenieril y pasa a formar parte del contenido cognitivo de la investigacion.

Por ello, la tesis se coloca deliberadamente en una frontera interdisciplinar. No pertenece unicamente al analisis de sistemas heterogeneos ni exclusivamente al aprendizaje automatico, sino al espacio en que ambos dominios se exigen mutuamente coherencia. Esta posicion hace mas exigente la argumentacion del manuscrito, pero tambien la vuelve mas valiosa: permite que cada resultado pueda leerse a la vez como proposicion sobre entrenamiento profundo y como proposicion sobre organizacion material del computo.

\section{1.3 Planteamiento formal del problema}
El entrenamiento de redes neuronales profundas sigue una secuencia bien definida de operaciones distribuidas entre las fases de propagación hacia adelante (forward) y propagación hacia atrás (backward). Durante el backward, es imprescindible conservar las activaciones intermedias generadas por cada capa, ya que son necesarias para el cálculo de gradientes mediante la regla de la cadena. Este requisito impone una demanda acumulativa y sostenida sobre el subsistema de memoria, particularmente en arquitecturas de redes neuronales que son entrenadas en Unidades de Procesamiento Gráfico (GPU).

En las implementaciones estándar, la totalidad del grafo computacional se asigna a un único dispositivo GPU, bajo la suposición de que su alto rendimiento en paralelo puede compensar las limitaciones físicas de capacidad de memoria. Sin embargo, la memoria disponible en la GPU es fija, jerárquicamente segmentada y no ampliable mediante abstracciones de software. Esta restricción se vuelve crítica al entrenar modelos profundos que procesan entradas de alta resolución o mantienen representaciones internas densas, como en redes tipo Transformers o en canalizaciones convolucionales 3D.

Las técnicas convencionales, incluyendo checkpointing, compresión de tensores, aritmética de baja precisión y rematerialización selectiva, alivian parcialmente la presión sobre la memoria, pero actúan exclusivamente dentro del dominio del acelerador. De forma análoga, los paradigmas de entrenamiento distribuido ofrecen escalabilidad en múltiples nodos GPU, pero no resuelven la saturación de memoria en configuraciones monodispositivo o en entornos con recursos locales limitados.

Por su parte, las Unidades Centrales de Procesamiento (CPU), que coexisten en la misma plataforma física, disponen de capacidades computacionales complementarias, espacios de memoria direccionables más amplios y soporte para ejecución vectorizada. A pesar de estas ventajas, permanecen subutilizadas durante el entrenamiento, no por insuficiencia arquitectónica, sino por la ausencia de un marco formal que permita integrarlas en la asignación de cargas bajo criterios consistentes de rendimiento y viabilidad.

Desde la perspectiva de la optimización discreta, este escenario constituye un problema de asignación por capas en un sistema heterogéneo, donde cada operación debe mapearse a un dispositivo de cómputo sujeto a restricciones observables: capacidad de memoria, latencia de transferencia, coste energético relativo y viabilidad de recomputación de activaciones, streaming o prefetching. El espacio de soluciones debe preservar la semántica de ejecución, respetar las dependencias del grafo y mantener la viabilidad global del sistema.

En consecuencia, la presente investigación aborda la siguiente brecha: la ausencia de un modelo de optimización formalmente definido, resoluble y parametrizado con métricas empíricas, que permita distribuir de forma eficiente las cargas de entrenamiento de los modelos de aprendizaje profundo entre CPU y GPU, minimizando la presión sobre la memoria de los aceleradores y preservando la viabilidad de ejecución, la eficiencia en el entrenamiento y la precisión del modelo en plataformas heterogéneas.

\section{1.4 Hipotesis de investigacion}
El entrenamiento supervisado en redes neuronales profundas implica la ejecución secuencial de operaciones tensoriales a través de múltiples etapas computacionales. Durante la fase de retropropagación, es necesario conservar en memoria las activaciones intermedias para calcular correctamente los gradientes mediante la regla de la cadena. Esta dependencia temporal genera una carga acumulativa sobre el subsistema de memoria de la GPU, cuya capacidad es físicamente limitada, jerárquicamente segmentada y no ampliable mediante software. Tales restricciones condicionan el volumen de activaciones que pueden almacenarse simultáneamente, comprometiendo la viabilidad del entrenamiento en modelos con alta densidad computacional o arquitecturas jerárquicas profundas, donde el almacenamiento intermedio se convierte en un recurso crítico.

En respuesta a este fenómeno, la hipótesis central de esta investigación plantea que es posible reducir la presión sobre la memoria de la GPU durante el entrenamiento de modelos de aprendizaje profundo mediante la distribución estructural de operaciones entre los subsistemas CPU y GPU, formalizada a través de un modelo de optimización discreta por capas, parametrizado con métricas reales de ejecución y aplicable a entornos heterogéneos, sin comprometer la exactitud del modelo ni la viabilidad de ejecución. Esta hipótesis implica que el entrenamiento de modelos de aprendizaje profundo puede reformularse como un problema de asignación computacional con restricciones, basado en recursos físicos medibles y verificables; que técnicas como la rematerialización y la transmisión interdispositivo pueden integrarse como decisiones internas dentro de un marco declarativo de optimización; y que el modelo resultante permite explorar configuraciones de ejecución no triviales, en las que la CPU actúa como un componente activo del grafo computacional, compensando las limitaciones de la GPU mediante la distribución inteligente de operaciones.

La validación de esta hipótesis se llevará a cabo mediante la construcción de un modelo de Programación Lineal Entera (ILP) que formalice la asignación óptima de operaciones, la simulación de su comportamiento sobre grafos de entrenamiento representativos y la comparación experimental de sus resultados frente a esquemas tradicionales centrados en GPU, evaluando indicadores como el uso máximo y promedio de memoria, la penalización en tiempo de ejecución, el consumo energético y la exactitud final del entrenamiento.

\section{1.5 Objetivos general y especificos}
El objetivo general de esta investigación es formular, implementar y validar empíricamente una estrategia estructurada de paralelización para el entrenamiento de redes neuronales profundas, que permita la distribución a nivel de capa de las cargas computacionales entre subsistemas CPU y GPU con el fin de reducir la presión sobre la memoria de los aceleradores y, simultáneamente, preservar la precisión del modelo, la viabilidad de ejecución y la eficiencia energética, bajo restricciones arquitectónicas realistas y métricas de rendimiento observables, aportando un marco metodológico reproducible para el análisis y la optimización del entrenamiento sobre entornos de cómputo heterogéneo.

Los objetivos específicos son los siguientes.

\begin{enumerate}
\item Caracterizar empíricamente el uso de memoria y el comportamiento de ejecución de arquitecturas de redes profundas representativas, identificando factores arquitectónicos y operativos que contribuyen a la saturación de la memoria de la GPU, mediante perfiles obtenidos en entornos reales de CPU y GPU.
\item Formalizar la asignación computacional por capas como un modelo de optimización discreta, incorporando variables de decisión para la colocación en los dispositivos, la persistencia de activaciones, las transferencias entre dispositivos y las estrategias de gestión de memoria, como la rematerialización, el streaming asíncrono y, cuando aplique, el prefetching, sujetas a restricciones de recursos y dependencias del grafo.
\item Parametrizar el modelo de optimización utilizando métricas observadas empíricamente, tales como huella de memoria por capa, latencia de ejecución, consumo energético y limitaciones específicas de cada dispositivo, garantizando que la formulación refleje dinámicas reales de entrenamiento.
\item Evaluar escenarios de ejecución híbrida CPU-GPU aplicando el modelo ILP a grafos de entrenamiento representativos, analizando el impacto de las decisiones de ubicación y recomputación sobre el uso de la memoria, el rendimiento temporal y la eficiencia de la asignación para el ajuste del modelo.
\item Validar experimentalmente el marco propuesto cuantificando el ahorro de memoria, las métricas de rendimiento, la eficiencia energética y la precisión final del modelo, e identificando las condiciones prácticas bajo las cuales la participación de la CPU mejora la eficiencia del entrenamiento frente a estrategias tradicionales centradas en la GPU.
\end{enumerate}

\section{1.6 Alcance decisional fijado (resultado de Fase 0)}
El alcance decisional del trabajo fue fijado explícitamente para evitar una deriva reduccionista de la contribución. La tesis no se compromete con una simple asignación estática única por capa, sino con una formulación fuerte que contempla asignación independiente para \textit{forward} y \textit{backward}, persistencia de activaciones como variable explícita del modelo y compatibilidad con estrategias de transferencia asíncrona y \textit{prefetching}. Esta decisión de alcance es metodológicamente importante porque impide confundir un prototipo parcial con la versión científica comprometida por la investigación.

La consecuencia epistemológica de este cierre es directa. El manuscrito puede utilizar formulaciones base y casos piloto como escalones explicativos, pero no puede presentar esas versiones intermedias como si agotaran la propuesta doctoral. Todo el texto debe leerse a la luz del contrato fuerte: la tesis aspira a una política de partición que no solo elija dispositivo por nodo, sino que gobierne también aspectos de disponibilidad de activaciones y sincronización entre subsistemas.

Desde esta perspectiva, el alcance no es un apartado administrativo, sino un mecanismo de disciplina científica. Fijar lo que la tesis realmente promete evita inflar aportes, ordena la validación por fases y mantiene coherencia entre escritura, software y evidencia. Esta coherencia es especialmente necesaria en un proyecto donde la complejidad técnica podría inducir a presentar como cierre doctoral lo que en realidad sería solo una capacidad exploratoria aún no sometida a contraste suficiente.

\section{1.7 Metodologia general y validacion}
La investigación se desarrolla a partir de un enfoque mixto que combina modelado formal y validación experimental, orientado a optimizar el entrenamiento de redes neuronales profundas bajo restricciones arquitectónicas verificables. El problema se formula de manera declarativa, representando las operaciones del grafo de entrenamiento como entidades discretas asignadas a CPU o GPU, sujetas a limitaciones de memoria, latencia, consumo energético y viabilidad estructural. Cada etapa metodológica se concibe como una instancia formalizable y verificable, sustentada en un modelo computacional explícito, mecanismos de trazabilidad algorítmica y criterios de validación reproducibles, garantizando consistencia y rigor científico.

El marco metodológico se estructura en cinco fases secuenciales e interdependientes. En la primera fase, se caracteriza empíricamente el comportamiento de cargas de trabajo de varios modelos de deep learning en plataformas físicas con CPU y GPU, identificando factores arquitectónicos y operativos que inciden en la saturación de memoria de la GPU. Para ello, se seleccionan arquitecturas representativas por su complejidad secuencial, densidad computacional y sensibilidad a la gestión de activaciones. Las ejecuciones reales se instrumentan con herramientas de perfilado y las métricas obtenidas, uso activo de memoria, duración de cómputo, tamaño de activaciones, latencias de transferencia, consumo energético y limitaciones específicas, se organizan en vectores temporales y matrices de observabilidad que constituyen la base parametrizable del modelo.

En la segunda fase, se formaliza un modelo declarativo sobre el grafo acíclico dirigido que describe el entrenamiento supervisado, donde cada nodo representa una operación diferenciable y las aristas codifican dependencias topológicas. El modelo se formula como un sistema de Programación Lineal Entera con variables binarias que definen la asignación de operaciones a CPU o GPU, la retención o recomputación de activaciones y el uso opcional de técnicas asíncronas como streaming interdispositivo, rematerialización o prefetching. Las restricciones incluyen límites de memoria por dispositivo, preservación del orden de ejecución, penalizaciones energéticas, latencias de transferencia y condiciones necesarias para el aprendizaje supervisado basado en gradientes.

La tercera fase consiste en parametrizar el modelo con las métricas reales obtenidas en la caracterización inicial, asegurando coherencia semántica entre observación y simulación, y minimizando divergencias entre ambos entornos. Se incluyen métricas de huella de memoria por capa, latencia de ejecución, consumo energético y limitaciones específicas. Este conjunto se complementa con un análisis de sensibilidad que varía sistemáticamente parámetros clave, penalizaciones por recomputación, latencia de transferencia y umbrales de memoria, para evaluar la estabilidad estructural de las soluciones y su capacidad de generalización.

Con el modelo parametrizado, en la cuarta fase se resuelven instancias mediante solvers ILP, derivando planes de ejecución que integran asignaciones CPU-GPU y políticas de gestión de memoria. A partir de cada solución, se implementa un simulador que reconstruye la secuencia de operaciones, verificando compatibilidad topológica, factibilidad técnica y cumplimiento de restricciones. El simulador estima el comportamiento agregado en términos de uso de recursos, penalizaciones temporales, solapamiento cómputo-comunicación y eficiencia de asignación. Estos resultados se validan de forma cruzada con ejecuciones físicas preliminares, ajustando discrepancias y comparando las estrategias híbridas con esquemas monolíticos centrados en GPU.

Finalmente, la quinta fase aborda la validación experimental en plataforma física, implementando las soluciones óptimas en PyTorch con asignación explícita de capas a CPU o GPU. Se miden uso activo de memoria, penalizaciones temporales, consumo energético estimado mediante telemetría interna y precisión final del modelo. Esta validación confirma tanto el rendimiento técnico como la coherencia metodológica del modelo, contrastando los resultados con esquemas tradicionales y verificando la solidez de sus supuestos en entornos heterogéneos reales.

En conjunto, este flujo metodológico integral y reproducible asegura que las decisiones de asignación optimizadas mantengan su viabilidad tanto en simulación controlada como en ejecución práctica, demostrando eficiencia operativa y consistencia de resultados bajo condiciones reales de cómputo heterogéneo.

\section{1.8 Contribucion cientifica}
La contribución científica de esta investigación radica en la formalización, validación empírica y aplicación de un modelo declarativo de optimización a nivel de capa para el entrenamiento de redes neuronales profundas en entornos heterogéneos CPU-GPU. Este modelo reinterpreta el entrenamiento supervisado como un problema estructurado de asignación, en el que las operaciones computacionales se redistribuyen de manera consistente, verificable y eficiente entre CPU y GPU, respetando las restricciones arquitectónicas observadas en plataformas reales.

A diferencia de los esquemas tradicionales centrados en la GPU, el enfoque propuesto busca reducir la presión de memoria activa sobre el acelerador mediante decisiones explícitas de rematerialización, desplazamiento computacional y sincronización entre dispositivos, preservando la exactitud del modelo y la viabilidad operativa. El sistema formal, basado en Programación Lineal Entera sobre grafos acíclicos de cómputo, integra métricas reales de ejecución tanto en la función objetivo como en las restricciones, garantizando que las soluciones generadas reflejen un comportamiento medible bajo condiciones reproducibles.

Metodológicamente, la investigación se apoya en la secuencia estructurada en la sección de metodología. Este proceso abarca la observación empírica, el modelado computacional, la resolución algorítmica, la simulación híbrida y la validación experimental. El flujo de trabajo está diseñado para ser replicable en distintos contextos arquitectónicos, diseños reales de redes neuronales y marcos de entrenamiento, y es extensible a sistemas multi-GPU, esquemas de memoria distribuida y arquitecturas heterogéneas emergentes con soporte para ejecución asíncrona en CPU.

Asimismo, el estudio introduce la noción operativa de sensibilidad arquitectónica, que evalúa la estabilidad de las soluciones óptimas frente a variaciones controladas en parámetros críticos, como latencia, penalizaciones energéticas y capacidad de retención de memoria. Este análisis no solo facilita la selección de configuraciones optimizadas, sino que también aporta información valiosa sobre el espacio de diseño técnico de los sistemas heterogéneos.

En consecuencia, la tesis presenta un marco computacional formalmente fundamentado y validado empíricamente para el entrenamiento de redes neuronales profundas con conciencia de recursos. Este marco permite razonar de forma sistemática sobre los compromisos entre memoria y rendimiento, apoya el desarrollo de estrategias de ejecución energéticamente eficientes y contribuye al rediseño de flujos de entrenamiento en entornos con saturación de GPU.

La originalidad de esta propuesta radica en integrar un modelo declarativo formal con métricas empíricas de ejecución para guiar la asignación óptima de operaciones en entornos heterogéneos. Su impacto potencial se extiende más allá del caso de estudio, ofreciendo a la comunidad de HPC y Deep Learning una herramienta metodológica replicable, adaptable a nuevas arquitecturas y capaz de influir en el diseño futuro de sistemas y algoritmos de entrenamiento distribuidos.

\section{1.9 Organizacion del documento}
La organización de la monografía responde al orden epistemológico del problema y no a una división meramente expositiva. El Capítulo 2 desarrolla los fundamentos teóricos que permiten comprender el entrenamiento profundo como un problema de ejecución sobre grafo, de jerarquía de memoria y de decisión bajo restricciones. El Capítulo 3 presenta la arquitectura metodológica del sistema como una cadena explícita de transformación de evidencia. El Capítulo 4 estudia el profiling empírico y la construcción de coeficientes observados. El Capítulo 5 formaliza la optimización ILP y sus mecanismos de robustificación, sensibilidad y comparación. El Capítulo 6 cierra el bucle entre planificación y ejecución mediante simulación y runtime híbrido. Finalmente, el Capítulo 7 sintetiza aportes, límites y proyección investigativa.

Esta secuencia es deliberada. Primero se fundamenta conceptualmente el problema; después se explica cómo se produce evidencia; luego se formula la decisión; y solo entonces se discute su ejecución real. En una tesis doctoral, esta progresión evita dos defectos comunes: teorizar sin anclaje experimental o describir resultados sin explicar con suficiente profundidad el marco conceptual que los hace significativos.
